% ECONOMICS_PROBLEM: {"id":"C73-71","title":"Non-equilibrium learning attractors","jel_code":"C73","source_id":"OP-0284"}
% FIRST_POSTED: 2026-10-10
% ENTRY_KIND: research_attempt
% ATTEMPT_STATUS: solved
\documentclass[11pt,a4paper]{article}
\usepackage[margin=25mm]{geometry}
\usepackage{amsmath,amssymb,amsthm}
\usepackage[hidelinks]{hyperref}
\hypersetup{pdftitle={Proof: A Sink without a Local Source That Is Not Attracting},pdfauthor={Ji Ho Bae}}
\setlength{\parindent}{0pt}
\setlength{\parskip}{4pt}
\setlength{\emergencystretch}{2em}
\allowdisplaybreaks[2]
\numberwithin{equation}{section}
\newtheorem{theorem}{Theorem}
\newtheorem{lemma}[theorem]{Lemma}
\newtheorem{proposition}[theorem]{Proposition}
\theoremstyle{definition}
\newtheorem{remark}[theorem]{Remark}
\newcommand{\R}{\mathbb R}
\newcommand{\Q}{\mathbb Q}
\DeclareMathOperator{\supp}{supp}
\DeclareMathOperator{\dist}{dist}
\title{Proof: A Sink without a Local Source\\
That Is Not Attracting}
\author{Ji Ho Bae\thanks{Prepared with GPT Codex assistance.}}
\date{10 October 2026}
\begin{document}
\maketitle
\begin{center}
\small\textbf{C73-71 \quad JEL C73 \quad Source OP-0284}\\
Proof with internal mathematical and source-fidelity review.
\end{center}
\begin{abstract}
We answer the local-source converse negatively for finite rational
bimatrix games. An explicit five-by-five game has a sink strongly
connected component consisting of every pure profile except one corner.
Its content has no local source in any product-face subgame, yet fails
uniform attraction under the replicator flow. A complete classification
of stationary support pairs in a four-by-four block excludes every
possible local source, including mixed sources involving the extra
strategies. The dynamical obstruction is a backward orbit approaching
a six-edge cycle. Exact logarithmic return coordinates give a linear
leading map with a strictly positive expanding eigenvector and a
uniformly bounded error. An invariant cone supplies the actual orbit
and controls it between sections. All payoff entries are rational;
no simulation or external switching approximation is needed.
\end{abstract}

\section{The question and the counterexample}\label{sec:question}

For a bimatrix game \(A,B\in\Q^{r\times s}\), let
\(X=\Delta_r\times\Delta_s\), where \(\Delta_j\) is the probability
simplex. Its replicator flow \(\varphi_t\) solves
\begin{equation}\label{eq:replicator}
 \dot x_i=x_i\bigl((Ay)_i-x^TAy\bigr),\qquad
 \dot y_j=y_j\bigl((x^TB)_j-x^TBy\bigr).
\end{equation}
The polynomial vector field gives a global flow on compact \(X\).
Each product face is invariant in both time directions, and a positive
coordinate stays positive at every finite time.

The preference graph has the pure profiles as vertices. A directed
edge is any unilateral weak payoff improvement between distinct
profiles; a payoff tie gives both directed edges. For a sink strongly
connected component \(H\), its content is
\[
 K_H=\{(x,y):\supp(x)\times\supp(y)\subseteq H\}.
\]
A compact invariant set \(K\) is uniformly attracting if some relative
positively invariant neighborhood \(U\) of \(K\) satisfies
\(\sup_{z\in U}\dist(\varphi_t(z),K)\to0\); this is the uniform
convention in~\cite[Definition 3.2]{BS}. The canonical C73-71 question
additionally calls a minimal such set an attractor~\cite{catalogue}.

A local source of \(H\) is a point \(z\in K_H\cap Y\), where \(Y\)
is a product-face subgame not contained in \(K_H\), which is a
quasi-strict Nash equilibrium of the \emph{negated} game on \(Y\).
Equivalently, all supported actions of each player have the same
original payoff, and every unsupported action available in \(Y\) has
strictly larger original payoff. There are no requirements on actions
absent from \(Y\). Biggar and Papadimitriou give the local-source
obstruction and ask its converse~\cite[Definition 2.1, Lemma 2.2,
Section 4]{BP}.

Set \(\eta=1/100\). Label the rows \(1,2,3,h,r\) and columns
\(1,2,3,d,c\), in that order, and define
\begin{equation}\label{eq:fullgame}
 A=\begin{pmatrix}
 2&0&-1&1&0\\
 -1&2&0&1&0\\
 0&-1&2&1&0\\
 3&3&3&0&0\\
 -\eta&-\eta&-\eta&-\eta&-1
 \end{pmatrix},\qquad
 B=\begin{pmatrix}
 2&-1&0&-1/2&-\eta\\
 0&2&-1&-1/2&-\eta\\
 -1&0&2&-1/2&-\eta\\
 1&1&1&2&-\eta\\
 0&0&0&0&-1
 \end{pmatrix}.
\end{equation}

\begin{theorem}\label{thm:main}
In \eqref{eq:fullgame}, let \(H\) be all 25 pure profiles except
\((r,c)\). Then \(H\) is one sink strongly connected component,
\[
 K_H=\{(x,y):x_r y_c=0\},
\]
and \(H\) has no local source in any product-face subgame. Nevertheless,
\(K_H\) is not uniformly attracting, and therefore is not an attractor.
This disproves the canonical local-source converse.
\end{theorem}

The proof establishes all three assertions for the same rational game.
The earlier catalogue attempts prove positive results for rectangles,
crosses and two-by-two games~\cite{prior1,prior2}; those results do not
settle this converse. Our six-edge set is not itself the sink component
of the full game, so the theorem on simple sink cycles in~\cite{cycles}
does not apply to it. We make no claim of an exhaustive literature search.

\section{The sink and every local-source candidate}\label{sec:sources}

Call rows \(1,2,3,h\) and columns \(1,2,3,d\) the resident block.
The old core on indices \(1,2,3\) has matrices
\[
 A_0=\begin{pmatrix}2&0&-1\\-1&2&0\\0&-1&2\end{pmatrix},
 \qquad B_0=A_0^T.
\]
A \emph{stationary support pair} has equal payoffs on each player's
positive support; unused actions are unrestricted. Thus it need not
be a Nash equilibrium of the full game.

\begin{lemma}\label{lem:sink}
The set \(H\) in Theorem~\ref{thm:main} is a sink strongly connected
component, and its content is \(\{x_r y_c=0\}\).
\end{lemma}
\begin{proof}
Read old indices cyclically modulo three. Write
\(F_i=(i,i+1)\) and \(G_i=(i,i-1)\); their payoffs are respectively
\((0,-1)\) and \((-1,0)\). The strict improvements
\(F_i\to G_i\to F_{i+1}\) form a six-cycle. Each old diagonal is
reachable from that cycle by an improvement to payoff 2. From a diagonal
there is a path
\[
 (i,i)\longrightarrow(h,i)\longrightarrow(h,d)
 \longrightarrow(j,d)\longrightarrow G_j,
\]
whose moving-player payoffs improve respectively from 2 to 3, from 1
to 2, from 0 to 1, and from \(-1/2\) to 0. These paths also contain
every remaining resident vertex, so the resident block is strongly
connected.

All \((r,j)\) with resident \(j\) communicate through column-payoff
ties. Each can move to \((h,j)\), whose row payoff is at least
\(0> -\eta\). Conversely, an old profile with row payoff \(-1\)
can move to row \(r\). All \((i,c)\) with resident \(i\) likewise
communicate through row-payoff ties; an old column payoff \(-1\)
allows entry, and a resident column payoff at least zero allows exit.
This proves strong connectivity of all of \(H\).

The only possible outgoing edges would move from an arm to \((r,c)\).
Both such moves lower the moving player's payoff from 0 to \(-1\),
so no weak outgoing edge exists. Finally a product distribution gives
positive mass to the sole omitted profile exactly when \(x_r y_c>0\).
The stated content follows; it is a compact invariant union of two faces.
\end{proof}

\begin{lemma}\label{lem:core}
Every stationary support pair of the old core is either pure, a
principal two-by-two pair with both payoffs \(4/5\), or the fully
mixed uniform pair with both payoffs \(1/3\).
\end{lemma}
\begin{proof}
Against a pure action, the other player's payoffs are the three
distinct numbers \(2,0,-1\); hence one-by-two and one-by-three
supports, and their reversals, are impossible. If one player uses all
three actions, indifference takes the form \(A_0z=v\mathbf1\).
Since \(\det A_0=7\) and \(A_0\mathbf1=\mathbf1\), normalization
forces \(z=(1/3,1/3,1/3)\) and \(v=1/3\). This excludes three-by-two
and two-by-three supports and identifies the full-support pair.

For two-by-two supports, cyclic relabelling reduces the row support
to \(\{1,2\}\). On column support \(\{1,3\}\), column 1 strictly
dominates column 3 for these rows. On column support \(\{2,3\}\),
row 2 strictly dominates row 1. Only the principal support remains.
There both indifference equations give masses \(2/5,3/5\) on
\(i,i+1\), and payoff \(4/5\). The pure profiles complete the list.
\end{proof}

\begin{lemma}[Resident sign certificate]\label{lem:sign}
Every stationary support pair of the resident four-by-four block has
at least one nonnegative payoff.
\end{lemma}
\begin{proof}
When neither \(h\) nor \(d\) is supported, Lemma~\ref{lem:core}
applies. Its pure payoffs are \((2,2)\), \((0,-1)\), and \((-1,0)\), so the claim
holds in that case.

Suppose \(d\) is supported and \(h\) is not. Column \(d\) earns
\(-1/2\). Two old columns cannot both tie it: the three coordinates
of \(A_0x\) sum to 1, so two coordinates equal to \(-1/2\) would
force the third to be 2. Attaining that maximum forces \(x\) to be
the corresponding pure action, whose other column payoffs are 0 and
\(-1\), a contradiction. If just one old column joins \(d\), all
old row payoffs are distinct: they are a positive multiple of
\(2,0,-1\), plus a common constant. Thus the row must be pure, but
then its old-column payoff cannot equal \(-1/2\). The only remaining
case is the pure column \(d\) with any old-row mixture, giving
payoffs \((1,-1/2)\).

If \(h\) is supported and \(d\) is not, row \(h\) earns 3, whereas
each old row earns at most 2. Therefore only \(h\) is played, with
payoffs \((3,1)\).

Finally suppose both hub actions are supported. If either old support
is empty, indifference forces both to be empty, giving \((h,d)\)
and payoffs \((0,2)\). Otherwise let \(\alpha,\beta\in(0,1)\)
be the old-column and old-row masses, and normalize their old mixtures
to \(\bar y,\bar x\). Indifference to the hubs gives
\[
 (A_0\bar y)_i=4-\frac1\alpha=:v,
 \qquad (A_0\bar x)_j=\frac1\beta-\frac32=:w
\]
on the respective old supports. Thus \((\bar x,\bar y)\) is one
of the core pairs in Lemma~\ref{lem:core}, and
\[
 \alpha=\frac1{4-v},\qquad \beta=\frac2{2w+3},\qquad
 (U_A,U_B)=\left(\frac3{4-v},\frac{4w+1}{2w+3}\right).
\]
The condition \(\beta<1\) requires \(w>-1/2\), excluding the
only core value \(w=-1\). All other core pairs have \(w\ge0\)
and \(v<3\), so these two resident payoffs are positive. This
exhausts the support cases.
\end{proof}

\begin{proposition}\label{prop:nosource}
No point of \(K_H\) is a local source in any product-face subgame.
\end{proposition}
\begin{proof}
A product-face subgame is not contained in \(K_H\) exactly when it includes
both row \(r\) and column \(c\). Consider any proposed source in
such a subgame. It must be stationary on its positive support.

If \(x_r=y_c=0\), it is a resident stationary pair. Quasi-strictness
in the negated game would require both resident payoffs to be strictly
less than \(-\eta\), because the added actions \(r,c\) each earn
\(-\eta\). Lemma~\ref{lem:sign} excludes this.

Suppose \(x_r>0,y_c=0\). If only row \(r\) is played, supported
columns earn 0, whereas column \(c\) earns \(-1\); hence the
negated-game Nash condition fails. Otherwise normalize the resident
row mass \(1-x_r\). Row indifference to \(r\) says that its resident
payoff is \(-\eta<0\). Since row \(r\) contributes zero to all
resident columns, their indifference equations give a resident
stationary pair. By Lemma~\ref{lem:sign} its column payoff \(v_B\)
is nonnegative. In the full profile, supported columns earn
\((1-x_r)v_B\ge0\), whereas adding \(c\) earns
\(-\eta(1-x_r)-x_r<0\). This also violates negated-game Nash.

If \(x_r=0,y_c>0\), the symmetric normalization argument uses the
resident column payoff \(-\eta\), and therefore a nonnegative
resident row payoff \(v_A\). Supported rows earn
\((1-y_c)v_A\ge0\), whereas adding \(r\) earns
\(-\eta(1-y_c)-y_c<0\). If only column \(c\) is played, the same
contradiction is directly \(0>-1\).
The remaining case \(x_r>0,y_c>0\) is outside \(K_H\). These cases
cover every support and every subgame, including omitted resident
actions; no condition on such omitted actions has been imposed.
\end{proof}

\section{An orbit approaching the cycle in negative time}\label{sec:orbit}

Set the hub probabilities to zero. Face invariance leaves the
four-by-four game on rows \(1,2,3,r\) and columns \(1,2,3,c\),
\begin{equation}\label{eq:facegame}
 A_* =\begin{pmatrix}
 2&0&-1&0\\-1&2&0&0\\0&-1&2&0\\
 -\eta&-\eta&-\eta&-1
 \end{pmatrix},\qquad B_*=A_*^T.
\end{equation}
Write \(\psi_t\) for the replicator flow with payoffs
\(-A_*,-B_*\). Negating both payoffs negates the vector field, so
\(\psi_t=\varphi_{-t}\) on this face.
Let \(S\) be the union of the six edges in the cyclic list
\begin{equation}\label{eq:cycle}
 (2,3),(1,3),(1,2),(3,2),(3,1),(2,1),(2,3).
\end{equation}
Every vertex and every edge of \(S\) is contained in \(K_H\).

\begin{lemma}\label{lem:orbit}
There is a point \(z_*\) with all eight probabilities positive in
\eqref{eq:facegame} such that
\(\dist(\psi_t(z_*),S)\to0\) as \(t\to\infty\).
\end{lemma}

The rest of this section proves the lemma directly from
\eqref{eq:replicator}. All constants used below are independent of
the scale of the logarithmic coordinates.

\subsection{Coordinates and the linear leading map}

Use an incoming section \(x_1=x_2\) with column 3 dominant. Define
\begin{equation}\label{eq:gaps}
 u=\log\frac{x_1}{x_3},\quad
 v=\log\frac{y_3}{y_2},\quad
 w=\log\frac{y_3}{y_1},\quad
 d=\log\frac{x_1}{x_r},\quad
 e=\log\frac{y_3}{y_c}.
\end{equation}
Every positive vector of these gaps determines unique strictly
positive probabilities on the section:
\[
 x_1=x_2=\frac1{2+\exp(-u)+\exp(-d)},\quad
 x_3=\exp(-u)x_1,\quad x_r=\exp(-d)x_1,
\]
\[
 y_3=\frac1{1+\exp(-v)+\exp(-w)+\exp(-e)},
\]
\[
 y_2=\exp(-v)y_3,\quad y_1=\exp(-w)y_3,\quad y_c=\exp(-e)y_3.
\]
For frozen play at \((1,3)\), the next column tie occurs at time
\(v\). At that time the unused gaps would be
\(u+3v,w+2v,d+(1-\eta)v,e-\eta v\), and
\(\log(x_1/x_2)=v\).
Exchange the players, then relabel the old indices simultaneously by
\(2\mapsto1,3\mapsto2,1\mapsto3\), exchanging \(r,c\) too.
Game \eqref{eq:facegame} is invariant under this operation. It turns
the new tie into the same incoming section. The resulting leading map is
\begin{equation}\label{eq:M}
 M\begin{pmatrix}u\\v\\w\\d\\e\end{pmatrix}
 =\begin{pmatrix}
 2v+w\\u+3v\\v\\e-\eta v\\d+(1-\eta)v
 \end{pmatrix}.
\end{equation}

Its resident block is
\[
 L=\begin{pmatrix}0&2&1\\1&3&0\\0&1&0\end{pmatrix},
 \qquad p(z)=\det(zI-L)=z^3-3z^2-2z-1.
\]
There is a unique positive root \(\lambda\in(3,4)\): the function
\(3/t+2/t^2+1/t^3\) strictly decreases from infinity to zero, and
\(p(3)<0<p(4)\). It is simple since \(p'(\lambda)>0\).
If \(z\) is any root, then
\(|z|^3\le3|z|^2+2|z|+1\), so \(|z|\le\lambda\).
Equality in the triangle inequality would force \(z^2,z,1\) to
have the same argument, and hence \(z=\lambda\). Thus every other
root has strictly smaller modulus. The remaining block of \(M\)
is the swap matrix, with eigenvalues \(1,-1\). Consequently
\(\lambda\) is a simple strictly dominant eigenvalue of \(M\).

Normalize the second coordinate to one. A corresponding eigenvector is
\begin{equation}\label{eq:eigenvector}
 q=\left(\lambda-3,1,\lambda^{-1},D,E\right),\qquad
 D=\frac{1-\eta(\lambda+1)}{\lambda^2-1},\qquad
 E=\lambda D+\eta.
\end{equation}
All coordinates are positive, and \(E-\eta=\lambda D>0\), because
\(\eta=1/100<1/(\lambda+1)\).

\subsection{A uniform error bound for the actual flow}

\begin{lemma}\label{lem:return}
On a sufficiently narrow closed cone about the ray \(\R_{>0}q\),
for all sufficiently large gap vectors \(z=(u,v,w,d,e)\), the first
subsequent tie \(y_2=y_3\) exists at time \(T=v+O(1)\). After the
relabeling above, the actual return map satisfies
\begin{equation}\label{eq:return}
 P(z)=Mz+O(1)
\end{equation}
uniformly on that cone.
\end{lemma}
\begin{proof}
By \eqref{eq:eigenvector}, choose the cone sufficiently narrow that,
for some fixed \(\mu>0\), the four affine functions
\begin{equation}\label{eq:ideal}
 u+3t,\quad w+2t,\quad d+(1-\eta)t,\quad e-\eta t
\end{equation}
are at least \(4\mu v\) for \(0\le t\le v\). The normalized
section \(v=1\) of the cone is compact. Thus all initial coordinates
are also at most a fixed multiple of \(v\).

Stop at the first of the following: the tie \(y_2=y_3\), the time
\(v+10\), or an unused gap \(u(t),w(t),d(t),e(t)\) reaching
\(2\mu v\). Before this stopping time, each unused probability
\(x_3,x_r,y_1,y_c\) is at most \(\exp(-2\mu v)\).
Put \(G(t)=\log(x_1/x_2)\) and
\(D_0(t)=\log(y_3/y_2)\). Exact payoff differences give
\begin{align}
 G'&=y_3+2y_2-3y_1\ge1-4y_1-y_c\ge\tfrac12,
       &G(0)&=0,\label{eq:G}\\
 D_0'&=-x_1+3x_2-2x_3
       =-1+4x_2-x_3+x_r,
       &D_0(0)&=v.\label{eq:D}
\end{align}
The last inequality in \eqref{eq:G} holds for large \(v\).
Therefore \(x_2\le e^{-t/2}\), and integration of \eqref{eq:D}
shows, throughout the stopped passage,
\begin{equation}\label{eq:timing}
 |D_0(t)-(v-t)|\le8+2(v+10)e^{-2\mu v}\le9
\end{equation}
for sufficiently large \(v\). There is no tie before \(v-9\).
If the other stopping conditions do not occur first, a tie must occur
by \(v+10\), since otherwise \eqref{eq:timing} would make its
strictly positive gap negative.

These estimates also bound the integrated deviations from the frozen
vertex, up to any such stopping time \(\tau\):
\begin{equation}\label{eq:integrals}
 \int_0^\tau x_2(t)\,dt\le2,\qquad
 \int_0^\tau y_2(t)\,dt\le
 e^9\int_0^{v+10}e^{-v+t}\,dt\le e^{19}.
\end{equation}
Here \(y_2\le e^{-D_0}\) was used. Each unused probability has
integral at most \((v+10)e^{-2\mu v}\).
The four exact unused-gap derivatives are
\begin{align}
 u'&=3-4y_2-5y_1-3y_c,\label{eq:unused}\\
 d'&=(1-\eta)-y_2-3y_1-(2-\eta)y_c,\notag\\
 w'&=2-x_2-5x_3-2x_r,\notag\\
 e'&=-\eta+x_2-2x_3+(\eta-1)x_r.\notag
\end{align}
By \eqref{eq:integrals}, their integrals differ from
\eqref{eq:ideal} by a fixed constant, independent of \(v\).
Up to \(v+10\), the ideal lower bound \(4\mu v\) decreases by
at most a fixed constant. For sufficiently large \(v\), the actual
gaps therefore remain strictly above \(2\mu v\). This rules out a
first bootstrap failure without assuming that the passage has already
completed.

The tie must consequently occur with \(v-9\le T\le v+10\).
Near that window, \eqref{eq:D} gives \(D_0'<-1/2\), so it is the
unique transverse crossing there. Equations \eqref{eq:G} and
\eqref{eq:integrals} yield \(G(T)=v+O(1)\), while
\eqref{eq:unused} yields the other four affine updates with bounded
errors. The relabelled gap vector is exactly
\((w(T),u(T),G(T),e(T),d(T))\), proving \eqref{eq:return}.
\end{proof}

\subsection{An invariant cone and the complete orbit}

Let \(\ell\) be a left \(\lambda\)-eigenvector of \(M\), normalized
by \(\ell(q)=1\). One may take a positive multiple of
\((1,\lambda,1/\lambda,0,0)\). Write
\(z=\tau q+h\), with \(\tau=\ell(z)\) and \(h\in\ker\ell\).
The restriction \(N=M|_{\ker\ell}\) has spectral radius strictly
less than \(\lambda\). Choose \(s\) between these two numbers.
There is a norm with \(\|Nh\|_*\le s\|h\|_*\); explicitly,
\[
 \|h\|_*:=\sum_{k=0}^{\infty}s^{-k}\|N^kh\|_\infty
\]
works. The series converges by the strict spectral-radius bound
(for example by Jordan form), and shifting its index proves the
stated inequality.

Choose \(\delta>0\) sufficiently small that
\[
 \mathcal C=\{\tau q+h:\tau>0,\ \ell(h)=0,
                         \ \|h\|_*\le\delta\tau\}
\]
lies in the cone of Lemma~\ref{lem:return}. All gap coordinates on
\(\mathcal C\) are positive and comparable with \(\tau\).
The bounded return error gives fixed constants \(C_1,C_2\) such that
\[
 \tau_{\rm new}\ge\lambda\tau-C_1,
 \qquad \|h_{\rm new}\|_*\le s\delta\tau+C_2.
\]
For sufficiently large \(\tau\),
\[
 (\lambda-s)\delta\tau>C_2+\delta C_1,
 \qquad \tau_{\rm new}\ge\frac{\lambda+1}{2}\tau.
\]
Thus the high-scale part of \(\mathcal C\) is invariant under the
actual return map, and the scales grow geometrically. A sufficiently
large multiple of \(q\) gives a nonempty initial section through
\eqref{eq:gaps}. Each successive passage exists by
Lemma~\ref{lem:return}; its duration is \(v_j+O(1)\), so the sum
of passage times tends to infinity.

It remains to control the trajectory between sections. The unused
probabilities tend uniformly to zero in each passage. For
\(t\le v_j/2\), \eqref{eq:timing} gives
\(y_2\le e^{-v_j/2+9}\); for \(t\ge v_j/2\),
\eqref{eq:G} gives \(x_2\le e^{-v_j/4}\).
In the first interval, setting the small unused probabilities and
\(y_2\) to zero and normalizing gives a point on the incoming edge
\(\{1,2\}\times\{3\}\). In the second interval, doing the same
with \(x_2\) gives a point on the outgoing edge
\(\{1\}\times\{3,2\}\). Normalization changes the state by a
quantity tending uniformly to zero. Both edges belong to \(S\).
The finite symmetry relabelings preserve \(S\); hence the same
estimate applies at every phase. Since \(v_j\to\infty\), the
entire trajectory approaches \(S\), proving Lemma~\ref{lem:orbit}.

\section{Completion and the obstruction}\label{sec:completion}

\begin{proof}[Proof of Theorem~\ref{thm:main}]
Lemma~\ref{lem:sink} identifies the sink and its content, and
Proposition~\ref{prop:nosource} excludes all local sources.
Choose the point \(z_*\) supplied by Lemma~\ref{lem:orbit}, and set
its two hub coordinates to zero in the full game. Its probabilities
on the four-by-four face are positive, so \(x_r y_c>0\) and
\(z_*\notin K_H\). By compactness,
\(d_*:=\dist(z_*,K_H)>0\).

Face invariance and exact time reversal give
\[
 \dist(\varphi_{-t}(z_*),K_H)\longrightarrow0.
\]
Every relative neighborhood \(U\) of compact \(K_H\) contains a
sufficiently thin metric neighborhood of \(K_H\). Therefore
\(\varphi_{-t}(z_*)\in U\) for all sufficiently large \(t\).
The flow identity then implies
\[
 \sup_{z\in U}\dist(\varphi_t(z),K_H)
 \ \ge\ \dist(\varphi_t(\varphi_{-t}(z_*)),K_H)
 \ =d_*>0.
\]
Uniform attraction fails for every neighborhood, including every
positively invariant one. This proves the theorem.
\end{proof}

The obstruction is an orbit that is asymptotic to a resident cycle in
backward time while both extra strategies remain positive at finite
times. Their favorable and unfavorable phases have geometrically
different durations. The positive eigenvector in
\eqref{eq:eigenvector}, together with the actual-flow error bound,
controls those durations and both extra coordinates. No stationary
point is used as a local source. The hub actions make all but one
corner a single sink without changing this invariant-face orbit.

\begin{remark}[Scope]
This is a counterexample to the entire main bimatrix implication,
under its weak-edge, all-subgame and uniform-attraction conventions.
It also refutes the corresponding unrestricted many-player implication
by adding players with one action. It does not establish failure of
the separately stated, weaker condition that every full-interior
initial state converges pointwise to \(K_H\) in forward time: the
present witness lies on a product face and uses time-shifted initial
states. An orbit starting in invariant \(K_H\) never leaves it.
\end{remark}

\begin{thebibliography}{9}\raggedright
\bibitem{BP} O. Biggar and C. H. Papadimitriou.
\emph{Sink equilibria and the attractors of learning in games}.
arXiv:2502.07975v4, 4 March 2026.
Definition 2.1, Lemma 2.2, and Section 4, Algorithmic problems, item 1.
\url{https://arxiv.org/html/2502.07975v4}.

\bibitem{BS} O. Biggar and I. Shames.
\emph{The Replicator Dynamic, Chain Components and the Response Graph}.
Proceedings of Machine Learning Research 201 (2023), 237--258.
Definition 3.2.
\url{https://proceedings.mlr.press/v201/biggar23a/biggar23a.pdf}.

\bibitem{cycles} O. Biggar and C. H. Papadimitriou.
\emph{Computing stable limit cycles of learning in games}.
arXiv:2602.11315v1, 11 February 2026. Theorem 4.1.
\url{https://arxiv.org/html/2602.11315v1}.

\bibitem{catalogue} \emph{Economics C73-71: Non-equilibrium learning
attractors}. Canonical statement, repository commit
\texttt{00ced55ba1223fb48603e1f15028d753997ee515}.
\url{https://github.com/mehmetmars7/Erdosproblems-llm-hunter/blob/00ced55ba1223fb48603e1f15028d753997ee515/attacks/open_problems/economics/statements/C73-71.tex}.

\bibitem{prior1} \emph{C73-71}, GPT 6 Astra Ultra, first catalogue attempt.
Same repository commit as above.
\url{https://github.com/mehmetmars7/Erdosproblems-llm-hunter/blob/00ced55ba1223fb48603e1f15028d753997ee515/attacks/open_problems/economics/gpt_6_astra_ultra/C73-71.tex}.

\bibitem{prior2} \emph{C73-71}, GPT 6 Astra Ultra, second catalogue attempt.
Same repository commit as above.
\url{https://github.com/mehmetmars7/Erdosproblems-llm-hunter/blob/00ced55ba1223fb48603e1f15028d753997ee515/attacks/open_problems/economics/gpt_6_astra_ultra/C73-71_v2.tex}.
\end{thebibliography}
\end{document}
